
\def\PUDcodigo{EME124}

\def\PUDcodacad{04506.43}

\def\PUDdisciplina{Refrigeração Industrial}

\def\PUDsemestre{S9}

\def\PUDhoras{80}

%NOVOS ---------------------------------------------------
\def\PUDhorasTeoria{76}
\def\PUDhorasPratica{4}

\def\PUDinterECA{\hyperref[PUD:EME115]{Mecanismos (04506.25)}
    
    \hyperref[PUD:EME114]{Termodinâmica (04506.27)}
    
    \hyperref[PUD:EME118]{Mecânica dos Fluidos (04506.34)}
    
    \hyperref[PUD:EME119]{Sistemas Mecânicos (04506.35)}
}

\def\PUDatuaisECA{---}

\def\PUDrecursos{Projetor multimídia; Quadro branco e pincel; Aulas em laboratório.}

%----------------------------------------------------------

\def\PUDcreditos{4}

\def\PUDprerequisito{ \hyperref[PUD:EME118]{EME118} }

\def\PUDpreacad{ \hyperref[PUD:EME118]{Mecânica dos Fluidos (04506.34)} }

\def\PUDobjetivos{Conhecer o princípio de funcionamento dos principais tipos de equipamentos de refrigeração encontrados na indústria: compressores, trocadores de calor, ventiladores, bombas, tubos, dutos e controles; Ser capaz de fazer uso efetivo da teoria da refrigeração industrial na prática da engenharia; Compreender os fenômenos relativos à refrigeração de alguma substância ou meio.}

\def\PUDementa{Refrigeração Industrial; Sistemas de Múltiplos Estágios de Pressão; Compressores Alternativos; Compressores Parafuso; Evaporadores, Serpentinas e Resfriadores; Recirculação de Líquido; Condensadores; Tubulações; Válvulas.}

\def\PUDconteudo{ & Refrigeração Industrial: Refrigeração industrial comparada ao condicionamento de ar para conforto;  O que é a refrigeração industrial;  Armazenamento de alimentos não congelados;  Alimentos congelados;  Processamento de alimentos;  Condicionamento de ar na indústria;  Refrigeração na indústria de manufatura;  Refrigeração na indústria da construção;  Refrigeração na indústria química e de processos;  (2 horas-aulas). 
\\ 
 & Carga Térmica: Qualidade do ar; Estimativa das trocas térmicas; Condições de projeto; Transmissão Térmica; Cargas de Infiltração e Ventilação; Componentes de carga térmica de resfriamento; Carga resultante da geração interna de calor; Carga térmica de insolação através de superfícies transparentes; Carga de insolação em superfícies opacas. 
 %Fundamentos da Termodinâmica Aplicada à Refrigeração: O Sistema Internacional de Unidades, SI;  Grandezas fundamentais e derivadas no Sistema Internacional;  Conversão de unidades;  Algumas constantes importantes do SI;  O diagrama pressão-entalpia;  A utilização das tabelas e dos diagramas de propriedades termodinâmicas dos refrigerantes;  O ciclo de refrigeração de Carnot;  O ciclo de Carnot com um refrigerante real;  O coeficiente de eficácia (COP);  Condições para COP elevados em ciclos de Carnot;  A bomba de calor de Carnot;  A equação da Conservação da Energia para Regime Permanente;  Análise do ciclo de Carnot através das entalpias;  A tonelada de Refrigeração;  Compressão de vapor seco comparada à compressão de vapor úmido;  O motor térmico comparado a um dispositivo de expansão;  O Ciclo Padrão de Compressão a Vapor e suas variantes;  Conclusão;  (4 horas-aulas). 
\\ 
 & Sistemas de Múltiplos Estágios de Pressão: Compressão em múltiplos estágios de pressão na Refrigeração Industrial;  A remoção do gás de "flash";  Resfriamento intermediário em compressão de duplo estágio;  Compressão com duplo estágio e uma única temperatura de evaporação;  A pressão intermediária ótima;  Compressão com duplo estágio e dois níveis de temperatura de evaporação;  Seleção do compressor;  Estágio único ou estágio duplo de compressão?;  Sistemas em cascata;  Conclusão;  (8 horas-aulas). 
\\ 
 & Compressores Alternativos: Tipos de compressores;  Rendimento volumétrico de espaço nocivo;  O efeito da temperatura de evaporação sobre a vazão de refrigerante;  O efeito da temperatura de evaporação sobre a capacidade frigorífica;  O efeito da temperatura de evaporação sobre a potência de compressão;  O efeito da temperatura de condensação sobre a vazão de refrigerante e a capacidade de refrigeração;  O efeito da temperatura de condensação sobre a potência de compressão;  Catálogos de fabricantes;  O rendimento volumétrico real;  Eficiência de compressão adiabática;  O efeito das temperaturas de evaporação e condensação sobre o COP;  Relação entre pressões e diferenças máximas de pressão;  O efeito do superaquecimento do vapor de aspiração e do sub-resfriamento do líquido;  Temperaturas de descarga e cabeçotes resfriados a água;  Lubrificação e resfriamento do óleo;  Controle da capacidade;  Compressores com múltiplas funções;  O mercado dos compressores alternativos;  (8 horas-aulas). 
\\ 
 %& Compressores Parafuso: Tipos de compressores rotativos parafuso;  Princípio de funcionamento;  Desempenho de um compressor de parafuso;  Eficiência de compressão adiabática;  O efeito das temperaturas de evaporação e de condensação;  Controle de capacidade e desempenho em carga parcial;  Compressores com relação entre volumes variáveis;  Injeção de óleo e resfriamento;  Aspiração a uma pressão intermediária;  Seleção do motor de acionamento;  O mercado dos compressores parafuso;  Compressor parafuso simples;  (6 horas-aulas). 
%\\ 
 & Evaporadores, Serpentinas E Resfriadores: Meios de transferência da carga de refrigeração;  O coeficiente Global de Transferência de Calor – U;  As aletas no lado do ar;  A mudança de fase do refrigerante no interior de tubos;  Propriedades do ar úmido - A Carta Psicométrica;  A lei da linha reta;  Linha do processo do ar numa serpentina;  O efeito de condições operacionais sobre o desempenho da serpentina;  Seleção de serpentinas em catálogos de fabricantes;  Controle da umidade em ambientes refrigerados;  Seleção e desempenho do ventilador e seu motor;  O número de serpentinas e sua localização;  Métodos de introdução do refrigerante e de controle de sua vazão;  Formação de neve em serpentinas de baixa temperatura;  Métodos de degelo de serpentinas;  Degelo por gás quente;  Serpentinas com borrifamento de anti-congelante;  Resfriadores de líquidos;  Temperatura ótima de evaporação;  (10 horas-aulas). 
\\ 
 & Recirculação de Líquido: O evaporador com recirculação de líquido;  Circulação por bombas e por pressão de gás;  Vantagens e desvantagens da recirculação de líquido;  Fundamentos da recirculação de líquido;  Admissão do refrigerante;  A recirculação por bomba;  Características das bombas de recirculação;  A recirculação de líquido por pressão de gás;  Análise energética do bombeamento por gás;  Considerações finais;  (8 horas-aulas). 
\\ 
 & Condensadores: Tipos utilizados na refrigeração industrial;  Condensação em superfícies exteriores;  A condensação no interior de tubos;  A Relação de Rejeição de Calor;  Desempenho de condensadores resfriados a ar e a água;  Torres de resfriamento;  Condensadores evaporativos;  Desempenho de condensadores evaporativos - características operacionais e de projeto;  O efeito da temperatura de bulbo úmido do ar ambiente;  O efeito das vazões do ar e da água sobre a capacidade;  Análise das condições favoráveis para a redução da vazão de ar;  Operação dos condensadores evaporativos durante o inverno;  Remoção de incondensáveis;  Tubulação em condensadores isolados;  Tubulação em condensadores paralelos;  O condensador evaporativo como meio de resfriamento para cargas exteriores ao ciclo frigorífico;  Tratamento da água em condensadores evaporativos;  O condensador como componente do ciclo frigorífico;  (8 horas-aulas). 
\\ 
 & Tubulações: Considerações gerais;  As funções das linhas de refrigerante;  Perda de carga em tubos de seção circular;  O diâmetro ótimo;  Dimensionamento da tubulação;  Linhas de líquido com trechos verticais;  Linhas horizontais e em elevação para misturas bifásicas;  Trechos em elevação na linha de aspiração de sistemas com expansão direta de refrigerantes halogenados;  (6 horas-aulas). 
\\ 
 & Válvulas: Tipos de válvulas;  Válvulas de bloqueio de atuação manual;  Válvulas de expansão manuais ou válvulas de balanceamento;  Válvulas de retenção;  Válvulas de solenóide;  Válvulas de solenóide pilotadas e acionadas por pressão de gás;  Válvulas reguladores de pressão: de ação direta, pilotadas e de compensação externa;  Controles de nível;  Válvulas de expansão controladas por superaquecimento;  Considerações finais;  (6 horas-aulas). }

\def\PUDmetodo{Aulas expositórias que podem ser teóricas e/ou práticas, onde as práticas no laboratório serão marcadas no decorrer da disciplina. A metodologia também contemplará o acompanhamento de um projeto real de um sistema de refrigeração existente no mercado.}

\def\PUDavalia{A avaliação será feita com aplicação de provas teóricas e/ou práticas, além da possibilidade de inclusão de trabalhos e seminários no decorrer da disciplina. Uma das avaliações será do projeto de uma sistema real com o devido memorial de cálculo.
}

\def\PUDbibbasica{ & \href{www.biblioteca.ifce.edu.br/index.asp?codigo_sophia=23840}{Stoecker}, W. F.  Refrigeração Industrial.  2ª ed. São Paulo, SP: Blucher, 2002.  371p.  ISBN: 8521203055. 
\\ 
 & \href{www.biblioteca.ifce.edu.br/index.asp?codigo_sophia=23691}{Macintyre}, A.  J.  Ventilação industrial e controle da poluição.  2ª ed.  Rio de Janeiro, RJ: LTC, 1990.  403p.  ISBN: 9788521611233.
\\ 
 &  \href{www.biblioteca.ifce.edu.br/index.asp?codigo_sophia=24058}{Moran}, M. J.;  Shapiro, H. N.;   Muson, B. R.;   DeWitt, D. P.  Introdução à engenharia de sistemas térmicos: termodinâmica, mecânica dos fluidos e transferência de calor.  Rio de Janeiro, RJ: LTC, 2005.  604p.  ISBN: 9788521614463.
}

\def\PUDbibcomplementar{ & \href{www.biblioteca.ifce.edu.br/index.asp?codigo_sophia=63133}{Fox}, R. W.; MacDolnad, A. T.; Pritchard, P. J.  Introdução à Mecânica dos Fluidos.  8ª ed.  Rio de Janeiro, RJ: LTC, 2014.  871p.  ISBN: 9788521623021. 
\\ 
 &  \href{www.biblioteca.ifce.edu.br/index.asp?codigo_sophia=61984}{Incropera}, F. P.; Witt, D. P.  Fundamentos de Transferência de Calor e de Massa.  7ª ed.  Rio de Janeiro, RJ: LTC, 2014.  672p.  ISBN: 9788521625049. 
\\ 
 &  \href{www.biblioteca.ifce.edu.br/index.asp?codigo_sophia=23804}{Silva}, N. F.  Compressores alternativos industriais: teoria e prática.  Rio de Janeiro, RJ: Interciência, 2009.  419p.  ISBN: 9788571932159.
\\
 &  \href{www.biblioteca.ifce.edu.br/index.asp?codigo_sophia=62615}{Telles}, Pedro Carlos da Silva.  Tubulações industriais: materiais, projeto, montagem.  10º ed.  Rio de Janeiro, RJ: LTC, 2014.  253p.  ISBN: 9788521612896.
\\
 &  \href{www.biblioteca.ifce.edu.br/index.asp?codigo_sophia=24513.}{Macintyre}, Archibald Joseph.  Bombas e instalações de bombeamento.  2º ed.  Rio de Janeiro, RJ: LTC, 2008.  782p.  ISBN: 9788521610861. } 
%\\ 
 %&  Maliska, C. R.;  Transferência de Calor e Mecânica dos Fluidos Computacional, 2ª edição, Editora: LTC.  ISBN : 9788521613961. }

\def\PUDautor{Francisco Frederico dos Santos Matos}

\def\PUDdata{2014-04-23}

\def\PUDrevisao{2}

\def\PUDrevisor{Francisco Frederico dos Santos Matos}

\def\PUDdatarevisao{2019-05-15}

\PUDcorpo
